\subsection{Class-Interference GNN}
\label{sec:class_interference_gnn}

Learning a new category can degrade old categories unevenly. Let
$\mathcal{M}_c$ be the Hungarian-matched queries whose targets have category
$c$. We define its category-conditioned DETR loss as
\begin{equation}
\mathcal{L}_c(\theta)=\frac{1}{\max(1,|\mathcal{M}_c|)}
\sum_{q\in\mathcal{M}_c}\left(\mathcal{L}_{\mathrm{cls}}(q,c)+
\lambda_{\mathrm{box}}\mathcal{L}_{\mathrm{box}}(q)\right),\quad
\mathcal{L}_{\mathrm{box}}=\lambda_1\mathcal{L}_1+
\lambda_{\mathrm{giou}}\mathcal{L}_{\mathrm{giou}}.
\label{eq:category_detection_loss}
\end{equation}
where $\mathcal{L}_{\mathrm{cls}}$ is the classification loss for $c$,
and $\mathcal{L}_1$ and $\mathcal{L}_{\mathrm{giou}}$ are matched-box
losses. The category gradient and compressed node feature are
\begin{equation}
g_c=\operatorname{vec}\left(\nabla_{\Theta_{\mathrm{FFN}}}
\mathcal{L}_c(\theta_{t-1})\right),\qquad
x_c=\operatorname{Normalize}(\operatorname{Pool}(g_c)),
\label{eq:class_gradient}
\end{equation}
where $g_c$ retains parameter-coordinate information and $x_c$ is the
fixed-dimensional feature supplied to the GNN.

For a source category $u$, a temporary low-rank probe update
$\Delta\theta^{\mathrm{probe}}_u$ is fitted using only source-category
supervision and then discarded. Its effect on an old category $c$ is
\begin{equation}
h_{u\rightarrow c}=\left[\mathcal{L}_c(\theta_{t-1}+
\Delta\theta^{\mathrm{probe}}_u)-\mathcal{L}_c(\theta_{t-1})\right]_+,
\label{eq:empirical_harm}
\end{equation}
where $h_{u\rightarrow c}$ is the observed positive increase in the old
category loss and supervises a directed edge.

The GNN performs directed message passing:
\begin{equation}
h_i^{(\ell+1)}=h_i^{(\ell)}+\Phi_\ell\left(h_i^{(\ell)},
\sum_{j\ne i}a_{j\rightarrow i}^{(\ell)}\Psi_\ell([h_j^{(\ell)},
h_i^{(\ell)},h_j^{(\ell)}-h_i^{(\ell)}])\right),
\label{eq:gnn_message_passing}
\end{equation}
where $a_{j\rightarrow i}^{(\ell)}$ is a normalized directed attention
coefficient and $\Phi_\ell,\Psi_\ell$ are learnable update functions.
The final edge predictor is
\begin{equation}
a_{u\rightarrow c}=\Psi_{\mathrm{edge}}([h_u^{(L)},h_c^{(L)},
h_u^{(L)}-h_c^{(L)}]),\qquad p_{u\rightarrow c}=\sigma(a_{u\rightarrow c}),
\label{eq:edge_prediction}
\end{equation}
where $p_{u\rightarrow c}$ is the predicted probability that learning $u$
interferes with old category $c$.

The GNN is trained on historical probe artifacts with
\begin{equation}
\mathcal{L}_{\mathrm{GNN}}=\operatorname{SmoothL1}(p_{u\rightarrow c},
\widetilde h_{u\rightarrow c})+\lambda_r[m_r-a_{u\rightarrow c^+}
+a_{u\rightarrow c^-}]_+,
\label{eq:gnn_loss}
\end{equation}
where $\widetilde h$ is the row-normalized harm target and
$c^+,c^-$ satisfy $h_{u\rightarrow c^+}>h_{u\rightarrow c^-}$. The
regression term fits harm magnitude and the ranking term preserves the
within-source ordering. If $\phi$ denotes GNN parameters, calibration uses
$\phi\leftarrow\phi-\eta_g\nabla_\phi\mathcal{L}_{\mathrm{GNN}}$;
$\phi$ is then frozen during detector adaptation.

For multiple current categories, their effects are aggregated as
\begin{equation}
s_t(c)=\max_{u\in K_t}p_{u\rightarrow c},\qquad
\mathcal{N}_t=\operatorname{TopK}_{c\in\mathcal{K}_{t-1}}s_t(c),
\label{eq:interference_neighborhood}
\end{equation}
where $\mathcal{N}_t$ contains old categories vulnerable to at least one
current category.

\subsection{Graph-Guided Detector Adaptation}
\label{sec:graph_guided_detector_adaptation}

We insert low-rank adapters into the final decoder FFN layers and expose the
current-category classifier rows as trainable detector capacity. Let
$\mathcal{P}_t$ denote this parameter subset and $\Delta\theta_t$ its
optimized update:
\begin{equation}
\theta_t=\theta_{t-1}+\Delta\theta_t,\qquad
\operatorname{supp}(\Delta\theta_t)\subseteq\mathcal{P}_t.
\label{eq:graph_guided_update}
\end{equation}
where parameters outside $\mathcal{P}_t$ remain frozen during the
incremental update. The neighborhood determines replay and protection signals
applied to this shared adaptation capacity; the GNN is not used at inference.

\subsection{Graph-Conditioned Continual Objective}
\label{sec:graph_conditioned_objective}

Every old category receives a base replay quota, while categories in
$\mathcal{N}_t$ receive additional examples. A frozen teacher predicts
missing old-category foreground objects in the same current or replay images;
confidence, duplicate, and overlap filtering produce completed targets
$\widetilde Y_n=Y_n\cup\widehat Y_n^{\mathrm{old}}$. Let
$\mathcal{C}_t^{\mathrm{loc}}=K_t\cup\mathcal{N}_t$,
$\mathcal{M}_t$ be matched queries with labels in this set, and
$\delta_t=\operatorname{vec}(\Delta\theta_t)$. We use
\begin{equation}
\begin{aligned}
\mathcal{L}_{\mathrm{local}}&=\frac{1}{\max(1,|\mathcal{M}_t|)}
\sum_{(q,y)\in\mathcal{M}_t}[m-z_{q,y}+
\max_{c\in\mathcal{C}_t^{\mathrm{loc}}\setminus\{y\}}z_{q,c}]_+,\\
\mathcal{L}_{\mathrm{off}}&=\frac{1}{\max(1,d_{\mathrm{off}})}
\left\|U_{\mathrm{off}}^\top\delta_t\right\|_2^2,\\
\mathcal{L}_{\mathrm{total}}&=\mathcal{L}_{\mathrm{det}}+
\lambda_{\mathrm{local}}\mathcal{L}_{\mathrm{local}}+
\lambda_{\mathrm{off}}\mathcal{L}_{\mathrm{off}}.
\end{aligned}
\label{eq:total_objective}
\end{equation}
where $m$ is the local margin, $z_{q,c}$ is a query logit,
$U_{\mathrm{off}}$ is an orthonormal basis of non-neighborhood gradient
directions, and $d_{\mathrm{off}}$ is its retained dimension.
$\mathcal{L}_{\mathrm{det}}$ is the weighted DETR classification, L1, and
generalized-IoU loss on completed current and replay targets. The local term
separates the new category from predicted competitors; the projection term
penalizes update components important to categories outside $\mathcal{N}_t$.

\begin{algorithm}[t]
\caption{Training the Class-Interference GNN}
\label{alg:gnn_training}
\begin{algorithmic}[1]
\Require Calibration detector $f_{\theta}$; category data; GNN parameters
$\phi$; calibration epochs $E$; probe step size $\eta_p$; GNN learning rate
$\eta_g$
\Ensure Trained GNN $G_{\phi}$ and directed interference scores
\State Compute $G_c=\partial\mathcal{L}_c/\partial\Theta_{\mathrm{FFN}}$
for each calibration category $c$ using Eq.~\ref{eq:class_gradient}
\State Construct node features $x_c$ from $g_c=\operatorname{vec}(G_c)$
\For{each source category $u$}
  \State Fit a temporary probe update $\Delta\theta_u^{\mathrm{probe}}$
  with step size $\eta_p$ using source-category supervision
  \State Measure $h_{u\rightarrow c}$ for every target category $c\neq u$
  using Eq.~\ref{eq:empirical_harm}
  \State Discard $\Delta\theta_u^{\mathrm{probe}}$
\EndFor
\For{$e=1$ to $E$}
  \State Predict directed scores $p_{u\rightarrow c}$ using
  Eqs.~\ref{eq:gnn_message_passing}--\ref{eq:edge_prediction}
  \State Compute $\mathcal{L}_{\mathrm{GNN}}$ using Eq.~\ref{eq:gnn_loss}
  \State Update $\phi\leftarrow\phi-\eta_g\nabla_{\phi}\mathcal{L}_{\mathrm{GNN}}$
\EndFor
\State Freeze $G_{\phi}$ and return the directed scores
\end{algorithmic}
\end{algorithm}

The trained GNN is subsequently used to predict the new-to-old interference
neighborhood for each incremental task. Its parameters are not updated by the
detector objective $\mathcal{L}_{\mathrm{total}}.
